\section*{अध्याय \ref{ch:b1:eukaryotic-cell} --- सुकेंद्रकी कोशिका का कार्यात्मक संगठन}

\begin{solution}{exo:b1:eukaryotic-cell:1}
\omterm{def:b1:eukaryotic-cell:nucleus}{केंद्रक} (गुणसूत्र, अनुलेखन); खुरदरी ER (स्रावी तथा झिल्ली-प्रोटीन);
चिकनी ER (लिपिड, कैल्सियम, विषहरण); \omterm{def:b1:eukaryotic-cell:endomembrane}{गॉल्जी} (छँटाई,
रूपांतरण, बाँधना); \omterm{def:b1:eukaryotic-cell:endomembrane}{लयनकाय} (पाचन); \omterm{def:b1:eukaryotic-cell:peroxisome}{परऑक्सीकाय} (ऑक्सीकरण,
कैटालेज़); \omterm{def:b1:eukaryotic-cell:mitochondrion}{सूत्रकणिकाएँ} (श्वसन, ATP)।
\end{solution}

\begin{solution}{exo:b1:eukaryotic-cell:2}
खुरदरी ER से बँधे राइबोसोमों द्वारा संश्लेषित और उसकी गुहा में पिरोया
जाता है (\omterm{prop:b1:eukaryotic-cell:secretory}{संकेत अनुक्रम}); ER से पुटिकाएँ निकलती हैं और सिस
\omterm{def:b1:eukaryotic-cell:endomembrane}{गॉल्जी} से मिल जाती हैं; प्रोटीन चट्टे को पार करता है, रूपांतरित और छाँटा
जाता है; वह ट्रांस फलक से किसी संघनन-रसधानी में निकलता है जो पककर एक
स्रावी कणिका बन जाती है; कणिका किसी संकेत पर शीर्ष झिल्ली से मिलती है और
वाहिनी में खाली हो जाती है।
\end{solution}

\begin{solution}{exo:b1:eukaryotic-cell:3}
$7 + 32 = \qty{39}{\%}$। भीतरी झिल्ली क्रिस्टी में इसलिए मुड़ी है कि वह
श्वसन-शृंखला और ATP सिंथेज़ ढोती है, जिनका निर्गम उसके क्षेत्रफल के
अनुपात में है; बाहरी झिल्ली एक सादा आवरण है।
\end{solution}

\begin{solution}{exo:b1:eukaryotic-cell:4}
पादप: सेलुलोस की भित्ति, बड़ी \omterm{def:b1:eukaryotic-cell:plantcell}{रसधानी}, \omterm{def:b1:eukaryotic-cell:plastid}{लवक} (\omterm{def:b1:eukaryotic-cell:plantcell}{जीवद्रव्यतंतु} भी)।
जांतव: \omterm{def:b1:eukaryotic-cell:endomembrane}{लयनकाय}, \omterm{def:b1:eukaryotic-cell:cytoskeleton}{तारककाय} (और रेंगने तथा भक्षण करने की
क्षमता)।
\end{solution}

\begin{solution}{exo:b1:eukaryotic-cell:5}
ER \qty{3}{min} पर (तब तक घट भी रही), \omterm{def:b1:eukaryotic-cell:endomembrane}{गॉल्जी} लगभग \qty{15}{min} पर,
कणिकाएँ लगभग \qty{60}{min} पर। ER से कणिका तक पारगमन लगभग
\qty{45}{min}; आधा लगभग \qty{110}{min} पर स्रावित।
\end{solution}

\begin{solution}{exo:b1:eukaryotic-cell:6}
\omterm{def:b1:eukaryotic-cell:endomembrane}{गॉल्जी}, जिसे सूक्ष्मनलिकाओं पर डाइनीन तारककेंद्र के पास थामे रखता है,
टूटकर बिखर जाता है; पुटिका-यातायात धीमा होकर विसरण पर आ जाता है और
परिधि तक परिवहन रुक जाता है; तर्कु नहीं बन पाता और \omterm{def:b1:cell-unit-of-life:cell}{कोशिकाएँ} समसूत्री
विभाजन में रुक जाती हैं; और पक्ष्माभ, जिनका गूदा सूक्ष्मनलिकाएँ हैं,
रुक जाते हैं (पहले से बने पक्ष्माभ स्थायी हैं, पर वे नए नहीं बनते)।
\end{solution}

\begin{solution}{exo:b1:eukaryotic-cell:7}
दो झिल्लियाँ, जिनमें भीतरी बैक्टीरिया जैसी है; हिस्टोनों के बिना
वर्तुल DNA; प्रतिजैविकों के प्रति सुग्राही बैक्टीरिया-आकार के राइबोसोम;
विदलन द्वारा विभाजन; और ऐल्फ़ाप्रोटियोबैक्टीरिया के साथ जुड़ते जीन-अनुक्रम।
वह अकेली नहीं जी सकती क्योंकि उसके अधिकांश जीन \omterm{def:b1:eukaryotic-cell:nucleus}{केंद्रक} में चले गए हैं: वह
अपने अधिकांश प्रोटीन \omterm{def:b1:cell-unit-of-life:cell}{कोशिकाद्रव्य} से आयात करती है।
\end{solution}

\begin{solution}{exo:b1:eukaryotic-cell:8}
शर्कराएँ ER और \omterm{def:b1:eukaryotic-cell:endomembrane}{गॉल्जी} में गुहा वाली ओर जोड़ी जाती हैं। हर पुटिका-संगलन
पक्षधरता बचाए रखता है: किसी पुटिका का गुहा-फलक अगले कक्ष का गुहा-फलक बन
जाता है, और जब पुटिका जीवद्रव्य-झिल्ली से मिलती है तो उसकी गुहा बाहर की
ओर खुल जाती है, इसलिए शर्कराएँ बाहर की ओर मुँह करती हैं। जो गुहा वाला है
वह, संस्थितिक रूप से, बाहर है।
\end{solution}

\begin{solution}{exo:b1:eukaryotic-cell:9}
हर कणिका $\pi d^2 = \qty{3.14}{\micro m^2}$; प्रति दिन
\qty{6280}{\micro m^2}, यानी शीर्ष फलक का \qty{157}{} गुना।
\omterm{def:b1:cell-unit-of-life:cell}{कोशिका} को उतनी ही दर से अंतःकोशिकता द्वारा झिल्ली वापस लेनी पड़ती है,
इसलिए शीर्ष झिल्ली हर नौ मिनट में आवर्तित होती है और वापस ली गई झिल्ली
फिर से बरतने के लिए \omterm{def:b1:eukaryotic-cell:endomembrane}{गॉल्जी} को लौटा दी जाती है।
\end{solution}

\begin{solution}{exo:b1:eukaryotic-cell:10}
\qty{320}{} अवशेषों की शृंखला वह पूर्ववर्ती है जिसमें उसका \qty{20}{}
अवशेषों का \omterm{prop:b1:eukaryotic-cell:secretory}{संकेत अनुक्रम} है, जो आम तौर पर ER में घुसते ही काट दिया जाता है।
पुटिकाओं के साथ, शृंखला संश्लेषण के दौरान गुहा में घुसती है, संकेत
कट जाता है, और झिल्ली उत्पाद को प्रोटीन-अपघटनी से बचाती है।
संकेत-रहित उत्परिवर्ती को ER का तंत्र पहचानता ही नहीं: वह मुक्त
राइबोसोमों पर बनता है और पूरी लंबाई में \omterm{def:b1:cell-unit-of-life:cell}{कोशिकाद्रव्य} में रह जाता है, कभी
स्रावित नहीं होता।
\end{solution}

\begin{solution}{exo:b1:eukaryotic-cell:11}
अंतःकोशिकता से भीतर ली गई झिल्लियाँ और घिसे \omterm{def:b1:eukaryotic-cell:organelle}{अंगक} पाचन के लिए \omterm{def:b1:eukaryotic-cell:endomembrane}{लयनकायों} तक
पहुँचाए जाते हैं; उस एंजाइम के बिना उनमें का लिपिड तोड़ा नहीं जा सकता और
\omterm{def:b1:eukaryotic-cell:endomembrane}{लयनकायों} को भर देता है, जो तब तक फूलते हैं जब तक \omterm{def:b1:cell-unit-of-life:cell}{कोशिका} विफल न हो जाए।
सबसे तेज़ झिल्ली-आवर्तन वाली \omterm{def:b1:cell-unit-of-life:cell}{कोशिकाएँ} (तंत्रिकाकोशिकाएँ, जो सूत्रयुग्मी
पुटिकाएँ लगातार पुनर्चक्रित करती हैं; बृहद्भक्षकाणु) उसे सबसे तेज़ जमा
करती हैं। लयनकायी एंजाइम जो शर्करा-संकेत ढोते हैं वही संकेत ढोता कोई
इंजेक्ट किया एंजाइम \omterm{def:b1:cell-unit-of-life:cell}{कोशिका} के पृष्ठ पर ग्राहियों से बँध सकता है और
ग्राही-मध्यस्थ अंतःकोशिकता द्वारा \omterm{def:b1:eukaryotic-cell:endomembrane}{लयनकायों} में लिया जा सकता है --- यही
एंजाइम-प्रतिस्थापन चिकित्सा का सिद्धांत है।
\end{solution}

\begin{solution}{exo:b1:eukaryotic-cell:12}
अंतःसहजीविता \omterm{def:b1:eukaryotic-cell:mitochondrion}{सूत्रकणिकाओं} और लवकों की अच्छी व्याख्या करती है (दो
झिल्लियाँ, अपना जीनोम तथा अपने राइबोसोम, विदलन, बैक्टीरिया संबंधी), और इस
तरह \omterm{def:b1:cell-unit-of-life:prokeuk}{सुकेंद्रकी} के ऊर्जा-उपापचय की भी। वह \omterm{def:b1:eukaryotic-cell:nucleus}{केंद्रक}, ER तथा \omterm{def:b1:eukaryotic-cell:endomembrane}{गॉल्जी}, या
कोशिकाकंकाल की व्याख्या नहीं करती, जिनका कोई बैक्टीरियल समकक्ष नहीं है और
जो पोषी वंशक्रम के आविष्कार जान पड़ते हैं --- उसकी झिल्ली के भीतरी मोड़ और
नए प्रोटीन। ``संघ'' \omterm{def:b1:eukaryotic-cell:organelle}{अंगकों} के लिए ठीक बैठता है; ``किसी झिल्ली-तंत्र से बँधा''
उस बड़े, अनसमझाए भाग की ओर संकेत करता है: कोई ऐसी पोषी \omterm{def:b1:cell-unit-of-life:cell}{कोशिका} जो
बैक्टीरिया को भीतर लेने से पहले ही अपनी रचना में \omterm{def:b1:cell-unit-of-life:prokeuk}{सुकेंद्रकी} थी।
\end{solution}

\begin{solution}{pb:b1:eukaryotic-cell:1}
\textbf{1.} कोई छोटी स्पंद अणुओं के एक सँकरे समवर्ग को अंकित करती है
ताकि उनकी स्थिति एक ही समय अंकित करे; और अतिरिक्त ठंडा ऐमीनो अम्ल आगे का
समावेशन रोक देता है ताकि वह समवर्ग बाद के संश्लेषण से धुँधला न पड़े।
\textbf{2.} ER: \qty{3}{min} (सबसे पहला बिंदु); \omterm{def:b1:eukaryotic-cell:endomembrane}{गॉल्जी} और संघनन-रसधानियाँ: लगभग \qty{15}{min}; कणिकाएँ: लगभग \qty{60}{min}।
\textbf{3.} ER से \omterm{def:b1:eukaryotic-cell:endomembrane}{गॉल्जी} तक लगभग \qty{10}{min}; \omterm{def:b1:eukaryotic-cell:endomembrane}{गॉल्जी} से कणिका तक लगभग
\qty{45}{min}।
\textbf{4.} अधिकतर कणिकाओं में (\qty{70}{\%}), दसवाँ भाग \omterm{def:b1:eukaryotic-cell:endomembrane}{गॉल्जी} में,
दसवाँ भाग स्रावित भी हो चुका: संघटक रास्ता और सबसे पहले बनी कणिकाएँ कुछ
छोड़ देती हैं, जबकि अधिकांश किसी उद्दीपन की प्रतीक्षा करती हैं।
\textbf{5.} पुटिकाओं का निकलना और उनका परिवहन ऊर्जा तथा तरल झिल्लियाँ
माँगते हैं; ठंड में ER-से-गॉल्जी वाला चरण रुक जाता है और प्रोटीन उस रुकावट
से पहले जमा हो जाता है।
\textbf{6.} प्रति \omterm{def:b1:cell-unit-of-life:cell}{कोशिका} प्रति दिन $10\,\mathrm{g}/\num{8e10} = \qty{1.25e-10}{g} =
\qty{125}{pg}$।
\textbf{7.} $V = \frac{\pi}{6}(1)^3 = \qty{0.52}{\micro m^3} =
\qty{5.2e-13}{mL}$; प्रोटीन $\qty{1.05e-13}{g} = \qty{0.1}{pg}$।
\textbf{8.} प्रति दिन $125/0.1 = 1250$ कणिकाएँ, यानी प्रति मिनट $0.87$।
\textbf{9.} $V = 12^3 = \qty{1728}{\micro m^3} = \qty{1.73e-9}{cm^3}$;
द्रव्यमान \qty{1.8e-9}{g} = \qty{1800}{pg}। वह रोज़ अपने द्रव्यमान का
$125/1800 = 7\%$ स्रावित करती है (पूरी \omterm{def:b1:cell-unit-of-life:cell}{कोशिका} केवल \qty{20}{\%} प्रोटीन
है, इसलिए अपनी प्रोटीन-मात्रा का एक-तिहाई)।
\textbf{10.} मोलर द्रव्यमान $250\times 110 = \qty{27500}{g/mol}$;
$\qty{1.05e-13}{g}/27500\times\num{6e23} = \num{2.3e6}$ अणु।
\textbf{11.} $250/5 = \qty{50}{s}$।
\textbf{12.} प्रति दिन अणु $1250\times\num{2.3e6} = \num{2.9e9}$;
एक राइबोसोम प्रति दिन $86\,400/50 = 1728$ बनाता है; चाहिए:
$\num{1.7e6}$ राइबोसोम, यानी \omterm{def:b1:cell-unit-of-life:cell}{कोशिका} के राइबोसोमों का लगभग \qty{40}{\%}।
\textbf{13.} $\pi d^2 = \qty{3.14}{\micro m^2}$।
\textbf{14.} $1250\times 3.14 = \qty{3900}{\micro m^2}$।
\textbf{15.} $3900/144 = 27$ गुना: \omterm{def:b1:cell-unit-of-life:cell}{कोशिका} उतना ही क्षेत्रफल
अंतःकोशिकता से वापस लेती है और \omterm{def:b1:eukaryotic-cell:endomembrane}{गॉल्जी} को लौटा देती है।
\textbf{16.} $20\,000/3900 = 5$ दिन: यदि कुछ भी न लौटता तो ER केवल पाँच
दिन की कणिका-झिल्ली दे पाती; झिल्ली को लगातार पुनर्चक्रित होना ही पड़ता है।
\textbf{17.} हर सूक्ष्मांकुर $\pi\times 0.1\times 1 =
\qty{0.31}{\micro m^2}$; 2000 के लिए \qty{630}{\micro m^2}; गुणक $5.4$।
\textbf{18.} शृंखला के आरंभ पर: \omterm{prop:b1:eukaryotic-cell:secretory}{संकेत अनुक्रम}, यानी बीस
अधिकतर जलविरोधी अवशेष।
\textbf{19.} संकेत ने बढ़ती शृंखला को पुटिकाओं में भेजा,
जहाँ किसी संकेत-पेप्टिडेज़ ने उसे हटा दिया; झिल्ली प्रोटीन को
प्रोटीन-अपघटनी से ढाल देती है।
\textbf{20.} केवल संकेत वाले प्रोटीन ही ER में घुसते हैं; किसी
कोशिकाद्रव्यी प्रोटीन में वह नहीं होता और वह बाहर ही रह जाता है।
\textbf{21.} एंजाइम मुक्त राइबोसोमों पर बनता है और \omterm{def:b1:cell-unit-of-life:cell}{कोशिकाद्रव्य} में रह
जाता है; और कोई पाचक एंजाइम होने के कारण (निष्क्रिय पूर्ववर्ती के रूप में
वह कम हानि करता है, पर सक्रिय हो जाए तो) वह \omterm{def:b1:cell-unit-of-life:cell}{कोशिका} को भीतर से पचा डालता।
\textbf{22.} ट्रांस \omterm{def:b1:eukaryotic-cell:endomembrane}{गॉल्जी} में: लयनकायी एंजाइम एक विशिष्ट चिह्न ढोता
है (कोई फ़ॉस्फ़ेटित शर्करा) जिसे कोई ग्राही पहचानता है और उसे स्रावी
कणिकाओं के बजाय \omterm{def:b1:eukaryotic-cell:endomembrane}{लयनकायों} को जाती पुटिकाओं में बाँध देता है।
\textbf{23.} ER में (और \omterm{def:b1:eukaryotic-cell:endomembrane}{गॉल्जी} में): पुटिका बनना, सूक्ष्मनलिकाओं पर
प्रेरक-परिवहन और झिल्ली-संगलन सब ATP खर्च करते हैं; और प्रोटीन-संश्लेषण
स्वयं रुक जाता है, इसलिए स्पंद पूरी भी नहीं होती।
\textbf{24.} प्रति मिनट $0.87$ कणिकाओं पर \qty{45}{min} का पारगमन:
यानी पारगमन में लगभग \qty{40}{} कणिकाओं जितना प्रोटीन।
\textbf{25.} संश्लेषण से कणिका तक पारगमन लगभग \qty{45}{min}
(ER गिनें तो \qty{55}{min}), और प्रयोग में लगभग दो घंटे बाद स्राव
(जीवित ग्रंथि में चार घंटे या अधिक, क्योंकि कणिकाएँ किसी भोजन की प्रतीक्षा
करती हैं); प्रति मिनट \qty{0.9}{} कणिकाएँ छूटती हैं; और प्रति दिन \omterm{def:b1:cell-unit-of-life:cell}{कोशिका}
के द्रव्यमान का \qty{7}{\%}।
\end{solution}
